% ===========================================================================
% NOVA SUBSECÇÃO 2.3.3, a colocar a seguir a "MÉTODOS POR APRENDIZAGEM E
% AGRUPAMENTO" e antes da secção "MÉTRICAS DE AVALIAÇÃO".
% ===========================================================================

\subsection{MÉTODOS APRESENTADOS AO DESAFIO SOCCERNET}
\label{sec:equipa_desafio}

As duas subsecções anteriores percorreram a literatura publicada, que aborda a
atribuição de equipa como um problema de separar os jogadores em dois grupos a partir
da sua aparência. As soluções apresentadas ao desafio SoccerNet, por atuarem dentro de
um sistema completo e serem avaliadas pela métrica oficial, revelam que o problema tem
na verdade duas etapas distintas, e que a segunda tem sido tratada com muito menos
atenção do que a primeira. Salvo indicação em contrário, as descrições que se seguem
resumem os textos submetidos pelas próprias equipas aos relatórios oficiais (Cioppa
et al., 2024; Giancola et al., 2025).

\subsubsection*{A estrutura comum}

Golovkin et al. (2025) sintetizam o panorama de forma útil. A generalidade dos sistemas
começa por extrair, para cada jogador, uma representação da sua aparência, seja um
histograma de cor, seja um vetor de reidentificação treinado por aprendizagem
contrastiva. Essas representações são depois agrupadas por métodos não supervisionados,
tipicamente o \emph{k-means} ou o DBSCAN, em dois a cinco grupos, consoante se
pretenda separar apenas as duas equipas ou também os árbitros e os guarda-redes. A
alternativa a esta estrutura é a aprendizagem de representações discriminativas de
equipa, como a proposta por Istasse et al. (2019), que dispensa o agrupamento por
produzir diretamente vetores que aproximam colegas de equipa e afastam adversários.

A solução de referência do desafio inscreve-se na primeira família, uma vez que aplica
o \emph{k-means} às representações que o PRTreID extrai, modelo que prevê em conjunto a
função e a afiliação de equipa de cada jogador.

\subsubsection*{Variantes exploradas pelos participantes}

Na edição de 2024, a equipa vencedora substituiu o agrupamento genérico por um modelo
treinado para o efeito, uma OSNet ajustada a 111 equipamentos distintos, e passou a
procurar cinco grupos em vez de dois, a saber, as duas equipas, os árbitros e os dois
guarda-redes. O procedimento é revelador do cuidado que a tarefa exige. Para estimar os
grupos das equipas consideram-se apenas os jogadores a menos de 30 metros do centro do
campo, de modo a excluir os guarda-redes, que são depois identificados separadamente
dentro das áreas. Cada jogador é finalmente atribuído ao grupo de menor distância do
cosseno.

Na edição de 2025 predominam duas orientações. A primeira é a delegação da tarefa no
mesmo modelo multimodal que lê o número de camisola, como faz a solução vencedora, que
pede ao modelo a cor do equipamento e agrupa depois por essa cor. A segunda é o reforço
da informação de aparência, como na solução que segmenta a silhueta do jogador e dela
extrai um descritor com as cores do equipamento e as respetivas áreas, usado a seguir
para a reidentificação, a função e a equipa.

Merece registo, pelo que ensina, o percurso da equipa que ficou em último lugar em 2024.
Propôs resolver a função e a equipa globalmente para todo o vídeo, através de uma rede
siamesa, de um agrupamento por \emph{k-means} com vários valores de $k$ e de uma floresta
aleatória que escolhia o agrupamento mais plausível a partir de características de
contexto de jogo. O resultado ficou abaixo da solução de referência, e os autores
atribuem-no ao facto de terem treinado o modelo com posições da anotação e o terem
avaliado sobre deteções reais. É um exemplo claro de como a qualidade das etapas
anteriores condiciona tudo o que vem depois.

\subsubsection*{A atribuição do lado do campo}

Separar os jogadores em dois grupos resolve metade do problema. Como a métrica oficial
exige que cada jogador seja identificado como pertencendo à equipa da esquerda ou à da
direita, é ainda necessário decidir qual dos dois grupos ocupa cada lado, e é aqui que
as soluções mais divergem.

A abordagem mais comum consiste em comparar a posição horizontal média dos jogadores de
cada grupo, atribuindo o lado direito àquele cuja média está mais à direita. É o que faz
a solução vencedora de 2025, que assume explicitamente que os jogadores tendem a
permanecer do seu lado do campo.

A solução vencedora de 2024 identifica, contudo, uma fragilidade desta suposição.
Substitui a média única por uma votação, em que cada \emph{frame} vota no grupo que nela
surge mais à direita, e justifica a alteração afirmando que a média simples ao longo de
todo o clipe pode falhar em sequências curtas ou quando apenas parte do campo está
visível.

Uma terceira via recorre à posição do guarda-redes. A equipa classificada em quarto
lugar em 2025 atribui a função de guarda-redes ao jogador mais afastado dentro de cada
área e submete depois o resultado a uma otimização que penaliza configurações
impossíveis, como a existência de dois guarda-redes do mesmo lado. Importa deixar claro,
por honestidade de atribuição, que \textbf{o recurso ao guarda-redes como referência de lado
não é original deste trabalho}, uma vez que já se encontra nesta e, de forma implícita,
na solução vencedora de 2024, que deteta os guarda-redes separadamente.

\subsubsection*{O que falta, e que este trabalho procura suprir}

Duas lacunas atravessam todas estas soluções.

A primeira é que \textbf{nenhuma quantifica com que frequência a suposição da posição média
falha}. A vencedora de 2024 afirma que falha, sem indicar em quantos casos nem em que
condições. Sem essa medida, não é possível saber se o problema é marginal ou estrutural,
nem justificar o custo de mecanismos alternativos.

A segunda é que a etapa de atribuição do lado nunca é avaliada isoladamente. O seu efeito
dilui-se na métrica global, que agrega a localização, o \emph{tracking}, a função, o
número e a equipa, o que impede saber quanto vale cada critério.

O Capítulo~\ref{cap:resultados} responde a ambas. Mede, sobre os 160 vídeos das três
divisões do conjunto de dados, a frequência com que o critério da posição média se
inverte, e avalia separadamente o contributo de cada nível da cascata de decisão
descrita na Secção~\ref{sec:met_lado}.
